\par\Needspace{4\baselineskip}
\begin{center}
\hypertarget{native-038-article021}{}\pdfbookmark[2]{21.}{article-021}21.
\end{center}
In der HYPERBEL würden $\varphi$ und $E$ imaginäre Größen. Will man solche Größen vermeiden, so muss man an ihrer Stelle andere Hilfsgrößen einführen. Den Winkel, dessen Kosinus $=\frac{1}{e}$ ist, haben wir oben bereits mit $\psi$ bezeichnet und für den Radiusvektor den folgenden Ausdruck gefunden.
\[
=\frac{p}{2e\cos\frac{1}{2}(v-\psi)\cos\frac{1}{2}(v+\psi)}
\]
Die Faktoren im Nenner dieses Bruchs, $\cos\frac{1}{2}(v-\psi)$ und $\cos\frac{1}{2}(v+\psi)$, sind für $v=0$ gleich. Der zweite verschwindet beim größten positiven Wert von $v$, der erste dagegen beim negativsten Wert. Setzt man also $\frac{\cos\frac{1}{2}(v-\psi)}{\cos\frac{1}{2}(v+\psi)}=u$, so gilt im Perihel $u=1$. Der Quotient wächst unbegrenzt, wenn sich $v$ seinem Grenzwert $180^0-\psi$ nähert. Dagegen nimmt er immer weiter ab, wenn man annimmt, dass $v$ zu seinem \SourceLeaf{039}anderen Grenzwert $-(180^0-\psi)$ zurückgeht. Somit entsprechen entgegengesetzten Werten von $v$ reziproke Werte von $u$, also Werte, deren Logarithmen einander entgegengesetzt sind.

Der Quotient $u$ eignet sich bei der Hyperbel sehr gut als Hilfsgröße. Fast ebenso zweckmäßig lässt sich an seiner Stelle der Winkel verwenden, dessen Tangens $=\operatorname{tang}\frac{1}{2}v\, .\,\BareRoot{\frac{e-1}{e+1}}$ ist. In Analogie zur Ellipse bezeichnen wir diesen Winkel mit $\frac{1}{2}F$. So ergeben sich leicht die folgenden Beziehungen zwischen den Größen $v$, $r$, $u$ und $F$, wobei wir $a=-b$ setzen, sodass $b$ eine positive Größe ist.
\[
\mathrm{I.}\quad b=p\operatorname{cotang}\psi^2
\]
\[
\mathrm{II.}\quad r=\frac{p}{1+e\cos v}=\frac{p\cos\psi}{2\cos\frac{1}{2}(v-\psi)\cos\frac{1}{2}(v+\psi)}
\]
\[
\begin{aligned}
\mathrm{III.}\quad\operatorname{tang}\frac{1}{2}F&=\operatorname{tang}\frac{1}{2}v\, .\,\BareRoot{\frac{e-1}{e+1}}\\
&=\operatorname{tang}\frac{1}{2}v\operatorname{tang}\frac{1}{2}\psi=\frac{u-1}{u+1}
\end{aligned}
\]
\[
\begin{aligned}
\mathrm{IV.}\quad u&=\frac{\cos\frac{1}{2}(v-\psi)}{\cos\frac{1}{2}(v+\psi)}=\frac{1+\operatorname{tang}\frac{1}{2}F}{1-\operatorname{tang}\frac{1}{2}F}\\
&=\operatorname{tang}(45^0+\frac{1}{2}F)
\end{aligned}
\]
\[
\begin{aligned}
\mathrm{V.}\quad\frac{1}{\cos F}&=\frac{1}{2}\left(u+\frac{1}{u}\right)\\
&=\frac{1+\cos\psi\cos v}{2\cos\frac{1}{2}(v-\psi)\cos\frac{1}{2}(v+\psi)}=\frac{e+\cos v}{1+e\cos v}\cdot
\end{aligned}
\]

Subtrahiert man auf beiden Seiten von Gleichung V $1$, so erhält man
\[
\begin{aligned}
\mathrm{VI.}\quad\sin\frac{1}{2}v\, .\,\BareRoot{r}&=\sin\frac{1}{2}F\, .\,\BareRoot{\frac{p}{(e-1)\cos F}}\\
&=\sin\frac{1}{2}F\, .\,\BareRoot{\frac{(e+1)b}{\cos F}}\\
&=\frac{1}{2}(u-1)\BareRoot{\frac{p}{(e-1)u}}=\frac{1}{2}(u-1)\BareRoot{\frac{(e+1)b}{u}}\cdot
\end{aligned}
\]

Addiert man entsprechend auf beiden Seiten $1$, so erhält man
\[
\begin{aligned}
\mathrm{VII.}\quad\cos\frac{1}{2}v\, .\,\BareRoot{r}&=\cos\frac{1}{2}F\, .\,\BareRoot{\frac{p}{(e+1)\cos F}}\\
&=\cos\frac{1}{2}F\, .\,\BareRoot{\frac{(e-1)b}{\cos F}}\\
&=\frac{1}{2}(u+1)\BareRoot{\frac{p}{(e+1)u}}=\frac{1}{2}(u+1)\BareRoot{\frac{(e-1)b}{u}}\cdot
\end{aligned}
\]

Teilt man VI durch VII, so erhält man wieder III. Multipliziert man die beiden Gleichungen, so ergibt sich
\[
\begin{aligned}
\mathrm{VIII.}\quad r\sin v&=p\operatorname{cotang}\psi\operatorname{tang}F=b\operatorname{tang}\psi\operatorname{tang}F\\
&=\frac{1}{2}p\operatorname{cotang}\psi\, .\,\left(u-\frac{1}{u}\right)=\frac{1}{2}b\operatorname{tang}\psi\, .\,\left(u-\frac{1}{u}\right).
\end{aligned}
\]

Aus der Verbindung der Gleichungen II und V erhält man ferner leicht
\[
\mathrm{IX.}\quad r\cos v=b\left(e-\frac{1}{\cos F}\right)=\frac{1}{2}b\left(2e-u-\frac{1}{u}\right)
\]
\[
\mathrm{X.}\quad r=b\left(\frac{e}{\cos F}-1\right)=\frac{1}{2}b\left\{e\left(u+\frac{1}{u}\right)-2\right\}.
\]
